% theme: graphs_of_trigonometric_functions←修正して下さい。
\MajorQuestionLesson{三角関数のグラフと最大・最小}

\SetRowSpace{3mm}

%y=sin(\theta)のような基礎の三角関数に定義域が与えられていて、
%図中の複数の点の座標を答えさせる問題
%4題 sin cos tan それぞれ１つは登場させて
\TypeQuestionSingle{}{}

\TypeQuestionDouble{基本となる三角関数のグラフをもとに考え、次の関数のグラフを、指定された範囲にかきなさい。}{horizontal_transform}
% 図生成工程：各小問に、横軸を弧度法で目盛った座標平面を挿入する。
%問題数を多くして下さい。合計12題、係数を変えてsin cos tanで4問ずつ。平行移動はなしでOKです。
%周期も答えさせるようにして下さい。
\QQRow
  {$y=\sin2\theta\quad(-\pi\leqq\theta\leqq\pi)$}{}
  {$y=\cos\dfrac{\theta}{2}\quad(-2\pi\leqq\theta\leqq2\pi)$}{}
\QQRow
  {$y=\tan2\theta\quad\left(-\dfrac{\pi}{2}<\theta<\dfrac{\pi}{2}\right)$}{}
  {$y=\sin\dfrac{\theta}{3}\quad(-3\pi\leqq\theta\leqq3\pi)$}{}

\TypeQuestionDouble{基本となる三角関数のグラフを平行移動して、次の関数のグラフを、指定された範囲にかきなさい。}{translation}
% 図生成工程：各小問に、横軸を弧度法で目盛った座標平面を挿入する。
%問題数を多くして下さい。合計12題、係数を変えてsin cos tanで4問ずつ。平行移動はx軸方向だけアリに統一してください。
%周期も答えさせるようにして下さい。
\QQRow
  {$y=\sin\left(\theta-\dfrac{\pi}{3}\right)\quad(-\pi\leqq\theta\leqq2\pi)$}{}
  {$y=\cos\left(\theta+\dfrac{\pi}{4}\right)\quad(-2\pi\leqq\theta\leqq\pi)$}{}
\QQRow
  {$y=\tan\left(\theta-\dfrac{\pi}{6}\right)\quad\left(-\dfrac{\pi}{2}<\theta<\dfrac{\pi}{2}\right)$}{}
  {$y=\sin\left(\theta+\dfrac{\pi}{2}\right)-1\quad(-2\pi\leqq\theta\leqq2\pi)$}{}
  
%ページを切り替えて最大・最小の問題

%y=sin(\theta-\dfrac{1}{6}\pi)のような単純な三角関数に定義域が与えられていて、
%最大・最小を求める問題
%8題
\TypeQuestionSingle{}{}

%cos　や sin　をtなどで置き換えて解く問題。
%相互関係を使う必要はナシのやつ
%4題
\TypeQuestionSingle{}{}



%cos　や sin　をtなどで置き換えて解く問題。
%相互関係を使う必要があるやつ
%4題
\TypeQuestionSingle{}{}




